\chapter{Investigate and Settle Decisions}\label{ch:decisions}
\section{Build an evidence base that answers the question}

\begin{headline}
  Start from the source that owns the behavior, and keep rules, observations, inferences
  and proposals apart.
\end{headline}

\begin{table}[H]
  \small
  \begin{tabularx}{\textwidth}{L{0.25\textwidth}Y}
    \toprule
    For & Inspect \\
    \midrule
    An analysis mismatch & The relevant definitions, implementation, input identity and
    output evidence. \\
    A workflow change & Its canonical procedure and current use. \\
    \bottomrule
  \end{tabularx}
  \caption{Where an investigation starts.}
  \label{tab:decisions-start}
\end{table}

\begin{rulebox}[Search safety]
  Search only explicit safe roots. Inspect known files and narrow remote paths; never scan
  the whole workspace or a mounted remote filesystem recursively \cite[Remote checkouts and
    execution]{skill:remote}.
\end{rulebox}

\begin{table}[H]
  \small
  \begin{tabularx}{\textwidth}{L{0.2\textwidth}Y}
    \toprule
    Statement & Says \\
    \midrule
    Rule & What the owning source requires. \\
    Observation & What a check actually found. \\
    Inference & What explains the observations. \\
    Proposal & A change not yet accepted. \\
    \bottomrule
  \end{tabularx}
  \caption{Four kinds of statement. For example, a successful run in one client is
    an observation supporting a client preference; it
    does not show that an unattended notification mechanism works \cite[Source traceability
      and corroboration]{skill:report-production}.}
  \label{tab:decisions-statements}
\end{table}

When evidence disagrees, find the first disagreement and the authority of each side:
\begin{itemize}
  \item a session can show that a run used a configuration; it cannot redefine the analysis
        convention;
  \item a policy file can require validation; it cannot prove that a run passed;
  \item a conclusion that depends on both cites both \cite[Source traceability and
          corroboration]{skill:report-production}.
\end{itemize}

\section{Use the planning pass for material choices}

\begin{headline}
  Before a material change, settle the owner's decisions in dependency order and confirm
  them as one record.
\end{headline}

The grilling pass applies before workflow or policy changes, architecture or
cross-repository work, physics or analysis design, and substantial report revisions
(Figure~\ref{fig:grilling}) \cite[Establish the decision tree; Work the
  frontier]{skill:grilling}.

\begin{figure}[H]
  \centering
  \begin{tikzpicture}[node distance=5mm and 6mm]
    \node[work,text width=30mm] (facts) {Inspect the facts};
    \node[work,text width=30mm,right=of facts] (split) {Separate owner choices from
      facts};
    \node[work,text width=30mm,right=of split] (order) {Order choices by dependency};
    \node[open,text width=30mm,below=8mm of order] (ask) {Ask the choices whose
      prerequisites are known};
    \node[open,text width=30mm,left=of ask] (confirm) {Confirm the decision record};
    \node[work,text width=30mm,left=of confirm] (impl) {Implement};
    \draw[flow] (facts)--(split);
    \draw[flow] (split)--(order);
    \draw[flow] (order)--(ask);
    \draw[back] (ask.east) to[bend right=50] node[right,lbl]{next\\round} (order.east);
    \draw[flow] (ask)--(confirm);
    \draw[flow] (confirm)--(impl);
  \end{tikzpicture}
  \caption{The grilling pass. Ordinary naming and formatting follow repository conventions
    without another round.}
  \label{fig:grilling}
\end{figure}

For example, a request to make an analysis script easier to understand first needs to
know whether it must preserve the scientific definition or implement a new one: that
decides the validation oracle. Then come the owning layer, the bounded input slice and
the acceptance evidence \cite[Establish the scientific contract]{skill:coding}.

The decision record keeps:
\begin{itemize}
  \item settled scope, non-goals and assumptions, with their evidence;
  \item affected owners and the implementation order;
  \item checks, the recovery path and remaining blockers.
\end{itemize}
Confirm it before implementation where the gate applies. If the owner already specified
every material choice, record them directly; do not manufacture a new round \cite[Confirm
  the decision record; Work the frontier]{skill:grilling}.

\section{Carry authorization with the action}

\begin{headline}
  Confirming a plan establishes understanding. A consequential action also needs an exact
  authorization boundary.
\end{headline}

Table~\ref{tab:consequence} is a compact guide; the consequential-actions procedure has
the full rules.

\begin{table}[H]
  \small
  \begin{tabularx}{\textwidth}{L{0.25\textwidth}Y}
    \toprule
    Action & What must be established \\
    \midrule
    Scoped inspection & Relevant target and safe read boundary. \\
    \rowcolor{ctpSurface0!30}
    Local reversible work & Assigned writing boundary and task scope. \\
    Own branch and review inbox & Nothing beyond the task: push the task's own branch or
    deliver a ready artifact to the reader's review inbox, then report it. \\
    \rowcolor{ctpSurface0!30}
    Shared change & Exact intent, target, state, recovery, and verification. Includes
    push to a main or shared branch, remote checkout update, and external PDF delivery. \\
    Operational action & Explicit task authorization for the job, service, mount,
    or other operation; credential initialization remains user-owned where required. \\
    \rowcolor{ctpSurface0!30}
    Destructive action & Specific authorization and recovery evidence for the exact
    target; cleanup additionally binds the immutable plan and batch. \\
    \bottomrule
  \end{tabularx}
  \caption{Action classes and their prerequisites
    \cite[Preflight]{skill:actions}.}
  \label{tab:consequence}
\end{table}

\begin{itemize}
  \item Record who authorized what, from which instruction, for which target and scope.
  \item The atomic steps of an authorized action need no new permission question each.
  \item Recheck when the target, state, scope, consequence or instruction changes.
  \item A successful inspection or a convenient recovery path never supplies missing intent
        \cite[Preflight]{skill:actions}.
\end{itemize}

\section{Decide, record, report}

\begin{headline}
  Between the gates, an agent does reversible work without asking and says what it did
  and why.
  A blocking scientific choice follows Figure~\ref{fig:drr}.
\end{headline}

\begin{figure}[H]
  \centering
  \begin{tikzpicture}[node distance=6mm and 8mm]
    \node[check,text width=42mm] (q) {A scientific choice blocks progress\\(a definition,
      which comparison answers the question)};
    \node[check,text width=34mm,right=12mm of q] (cheap) {Can a cheap run settle it?};
    \node[work,text width=34mm,right=12mm of cheap] (run) {Run it and decide from the
      result};
    \node[check,text width=34mm,below=10mm of cheap] (def) {Is a default defensible, and
      the purpose unchanged?};
    \node[work,text width=34mm,right=12mm of def] (rec) {Take the default best supported by the
      sources and what is already agreed, record it, keep working};
    \node[open,text width=34mm,below=10mm of def] (ask) {Ask the user before acting};
    \draw[flow] (q)--(cheap);
    \draw[flow] (cheap)--node[above,lbl]{yes}(run);
    \draw[flow] (cheap)--node[right,lbl]{no}(def);
    \draw[flow] (def)--node[above,lbl]{yes}(rec);
    \draw[flow] (def)--node[right,lbl]{no}(ask);
  \end{tikzpicture}
  \caption{Decide, record, report \cites{skill:prototype}[Decide, record,
      report]{skill:actions}.}
  \label{fig:drr}
\end{figure}

A recorded default goes into the session's table of agent defaults with the alternative,
the reason, the evidence and how to reverse it. Every result that depends on an
unreviewed default stays awaiting human review, and the next report lists those defaults
first. The work keeps moving while the scientific judgment stays visible and reversible.

\begin{rulebox}[Never a default]
  A confirmed decision record or an agreed convention is never a default to override.
\end{rulebox}
